\chapter{Introduction}\label{ch:intro}

\section{The problem, in one paragraph}

A flying ad-hoc network of unmanned aerial vehicles must do two things at once that pull against
each other. It must broadcast telemetry often enough to be useful --- position, velocity, battery,
mode, tens of times a second --- and it must make that telemetry \emph{tamper-evident}, so that a
record cannot be forged, replayed or silently altered by a participant or an outsider. The second
requirement costs bytes. A digital signature at the 128-bit security level is $64$\,B for Ed25519
or ECDSA-P256, and $96$\,B for the BLS12-381 variant actually shipped by the reference
implementation. A telemetry record, once compressed, is $45$\,B.

That inequality is the whole thesis. \textbf{The signature is larger than the thing it signs.}

\section{Why this is not simply a compression problem}

The obvious response --- compress harder --- fails, and it fails for a reason worth stating
precisely. Payload compression was the subject of prior work; this thesis begins where that work
ends. Once a record is squeezed from $191$\,B of JSON to $45$\,B of delta-encoded CBOR, the
authentication object has gone from $25\%$ of the on-air cost to $59\%$. Every further byte saved
on the payload increases the fraction spent on cryptography. Compression does not solve the
problem; \emph{compression is what creates it}.

The design space that remains has four axes, and this thesis searches them jointly:

\begin{itemize}
\item \textbf{Encoding} $e$ --- how a record is serialised (JSON, CBOR, MessagePack, delta-CBOR).
\item \textbf{Placement} --- \emph{where} the signature sits relative to the records it covers:
      inline per record, self-batched over $b$ of one sender's records, cross-signer aggregated,
      or block-level across multiple frames.
\item \textbf{Scheme} $\sigma$ --- which signature algorithm, and therefore how many bytes the
      authentication object costs and how much CPU verification demands.
\item \textbf{Batch} $b$ --- how many records share one authentication object, bounded by the
      frame size on one side and by a freshness deadline on the other.
\end{itemize}

\section{The question this thesis actually asks}

The literature on efficient authentication asks \emph{how many bytes does my scheme save}. That is
a useful question and it has been answered many times. This thesis asks a different one, because
for the deployments that motivate it the byte count is not the binding constraint:

\begin{quote}
\textbf{Which configurations can run at all?}
\end{quote}

That reframing --- from optimisation to feasibility --- is the thesis's organising decision, and
\Cref{ch:reproducibility} records why it was made partway through the work rather than at the
start. Briefly: as a co-design optimisation the result is mid-tier, because the optimisation is
largely closed-form and a factorial ablation showed only one genuine interaction among the four
axes. As a \emph{boundary} result the same evidence is stronger, because an impossibility cannot
drift as models improve --- and four of this project's own performance numbers did drift, and had
to be corrected.

\section{Research questions}

\begin{description}
\item[RQ1] How large is authentication overhead after compression, and where is the payload-size
      threshold below which it dominates?
\item[RQ2] Which authentication placement minimises on-air bytes while keeping records verifiable
      under loss, and what does self-batching change relative to cross-signer aggregation?
\item[RQ3] When does an aggregatable scheme beat a fast conventional one, as a function of relay
      fraction, arrival rate and the radio/CPU power ratio?
\item[RQ4] What is the end-to-end optimal configuration, how much does it beat the naive and
      compression-only baselines, and do the analytical models survive contention simulation and
      real hardware?
\item[RQ5] \emph{(added during the work)} Where is the boundary beyond which no choice of
      \emph{existing} cryptography is feasible at all --- and how much of that boundary is
      arithmetic rather than an artifact of this design?
\end{description}

RQ5 was not in the original charter. It emerged from the low-rate arm and became the thesis's most
durable result; \Cref{ch:lowrate} develops it and is honest about which part of it survives
scrutiny and which part does not.

\section{Contributions}

\begin{enumerate}
\item \textbf{An overhead-dominance characterisation} for compressed authenticated FANET ledgers,
      with the threshold stated in closed form and measured across four encodings
      (\Cref{ch:theory}, \Cref{ch:bytes}).

\item \textbf{A placement theory} distinguishing self-batching from cross-signer aggregation, with
      a loss-robustness frontier proving that frame-level batching Pareto-dominates block-level
      under any stationary loss process --- not merely under the independent-loss assumption the
      emulator implements (\Cref{ch:theory}).

\item \textbf{An exclusion bound instantiated on a real band, separated by durability.} Three EU868
      data rates are excluded \emph{unconditionally}: a $64$\,B signature will not fit a $51$\,B
      payload at a zero-byte header. A fourth is excluded by \emph{our own} frame header, and the
      thesis quantifies that contingency rather than leaving it implicit (\Cref{ch:lowrate}).

\item \textbf{A feasibility envelope} rather than a byte ratio: the largest neighbourhood each
      configuration sustains, at $1.9\times$--$3.2\times$ the inline baseline, reported at two
      thresholds because they answer different questions (\Cref{ch:codesign}).

\item \textbf{Validated models with their gaps stated}, including a broadcast channel model
      validated to $\leq 0.08\%$ against an independent slot-exact simulator and $\leq 0.51\%$
      against NS-3, plus two-radio hardware measurement of the airtime constants
      (\Cref{ch:validation}).

\item \textbf{An ablation that narrows our own claim}, showing that only placement and batching
      genuinely interact, and that an apparent encoding interaction was a ratio-scale artifact
      (\Cref{ch:codesign}).

\item \textbf{A reproducible testbed and a defect taxonomy.} Every model-derived artifact is
      re-derived byte-identically by an automated gate, and the classes of defect that survived
      that gate are catalogued --- including the ones found in this work
      (\Cref{ch:reproducibility}).
\end{enumerate}

\section{Explicit non-goals}

This thesis introduces no new cryptographic primitive, no new consensus protocol and no
post-quantum implementation; post-quantum sizes appear only as analytical projections. Key
compromise, multi-hop routing research and finality semantics are out of scope. The low-rate arm is
a \emph{modelling} chapter: no LoRa radio was metered and no energy column exists for it.

\section{Structure of this thesis}

\Cref{ch:background} places the work in its literature. \Cref{ch:sysmodel} fixes the system and
threat models and the notation. \Cref{ch:theory} develops the theory. \Cref{ch:implementation}
describes what was built and \Cref{ch:methodology} how it was measured. \Cref{ch:bytes} and
\Cref{ch:codesign} give the 802.11 results; \Cref{ch:validation} the channel-model validation and
hardware; \Cref{ch:lowrate} the low-rate regime and the exclusion bound.
\Cref{ch:reproducibility} is unusual for a thesis and deliberate: it reports the errors found in
this project's own results, how they were found, and what class of defect resisted the safeguards.
\Cref{ch:conclusions} concludes.
